\documentclass[12pt,a4paper,twocolumn]{article}
 
% -------------------- PAQUETES --------------------
\usepackage[utf8]{inputenc}
\usepackage[T1]{fontenc}
\usepackage[english,spanish,es-tabla]{babel} 
\usepackage[a4paper,margin=2.5cm]{geometry}
\usepackage{newtxtext}
\usepackage{newtxmath}
\usepackage{setspace}
\usepackage{amsmath}
\usepackage{graphicx}
\usepackage{float}
\usepackage{booktabs}
\usepackage{array}
\usepackage{tabularx}
\usepackage{caption}
\usepackage[hidelinks]{hyperref}
\def\UrlBreaks{\do\/\do\-\do\.\do\=\do\_}
\Urlmuskip=0mu plus 1mu
 
% -------------------- FORMATO --------------------
\setstretch{1.15}
\setlength{\columnsep}{0.7cm}
\setlength{\parindent}{0.5cm}
\setlength{\parskip}{0pt}
\emergencystretch=1.5em
\captionsetup{font=small,labelfont=bf,labelsep=period,skip=4pt,justification=justified,singlelinecheck=true}
\setlength{\textfloatsep}{10pt plus 2pt minus 2pt}
\setlength{\dbltextfloatsep}{10pt plus 2pt minus 2pt}
\setlength{\intextsep}{8pt plus 2pt minus 2pt}
\renewcommand{\topfraction}{0.9}
\renewcommand{\dbltopfraction}{0.9}
\renewcommand{\textfraction}{0.1}
\renewcommand{\floatpagefraction}{0.8}
\renewcommand{\dblfloatpagefraction}{0.8}
\setcounter{topnumber}{3}
\setcounter{dbltopnumber}{3}
\newsavebox{\tabbox}
\newcommand{\fitcol}[1]{\sbox{\tabbox}{#1}\ifdim\wd\tabbox>\columnwidth\resizebox{\columnwidth}{!}{\usebox{\tabbox}}\else\usebox{\tabbox}\fi}
\newcommand{\fuente}[1]{\par\smallskip{\small\textit{Fuente:} #1\par}}
\setlength{\tabcolsep}{4pt}
\renewcommand{\arraystretch}{1.1}
 
\begin{document}
 
\twocolumn[{%
% -------------------- TITULO --------------------
\begin{center}
{\fontsize{16}{19}\selectfont
\textbf{Comparación experimental del rendimiento de
Bubble Sort y Merge Sort bajo distintos escenarios de entrada}
\par}
 
\vspace{0.3cm}
 
{\fontsize{14}{17}\selectfont
\textbf{Experimental Comparison of Bubble Sort and Merge Sort
Performance under Different Input Scenarios}
\par}
 
\vspace{0.5cm}
 
{\normalsize
Omar Alejandro Flores Anahua\par}
 
{\small
\textit{Escuela Profesional de Ingeniería en Informática y Sistemas,}\\
\textit{Universidad Nacional Jorge Basadre Grohmann, Tacna, Perú}\\
\texttt{ofloresa@unjbg.edu.pe}
\par}
\end{center}
 
\vspace{0.3cm}
 
% -------------------- RESUMEN --------------------
\begin{otherlanguage}{spanish}
{\fontsize{11}{13}\selectfont
 
\noindent\textbf{Resumen:}
\textit{\textbf Comparar experimentalmente el rendimiento de
Bubble Sort y Merge Sort, implementados en C++, bajo distintos tamaños
y patrones de entrada. Se programaron ambos algoritmos
sin funciones de ordenación de la biblioteca estándar y se evaluaron cinco
patrones de entrada (aleatorio, ascendente, descendente, casi ordenado y con
duplicados) con $n = 1\,000$, $10\,000$ y $100\,000$, con 30 repeticiones por
condición y semillas fijas; Bubble Sort no se midió con $n = 100\,000$ por su
costo de ejecución. Se registraron el tiempo, las comparaciones, los
intercambios o movimientos y la profundidad de recursión. Con $n = 10\,000$ y datos aleatorios, la mediana del tiempo fue de 285,4 ms en Bubble Sort y de 1,75 ms en Merge Sort, unas 163 veces
menor. Bubble Sort realizó siempre $n(n-1)/2$ comparaciones y su tiempo varió
con los intercambios; Merge Sort tardó entre 10 y 20 ms con $n = 100\,000$,
según el patrón. El crecimiento medido fue compatible con $O(n^2)$ y $O(n\log n)$, y el patrón de entrada modificó los tiempos en un factor cercano a 2 en ambos algoritmos. La ventaja de Merge Sort se obtuvo a cambio de memoria auxiliar de orden $O(n)$.}
 
\vspace{0.15cm}
 
\noindent\textbf{Palabras clave:}
algoritmos de ordenamiento, Bubble Sort, Merge Sort,
rendimiento, análisis experimental.
 
}
\end{otherlanguage}
 
\vspace{0.3cm}
 
% -------------------- ABSTRACT --------------------
\begin{otherlanguage}{english}
{\fontsize{11}{13}\selectfont
 
\noindent\textbf{Abstract:}
\textit{\textbf To experimentally compare the performance of
Bubble Sort and Merge Sort, implemented in C++, under different input sizes
and data patterns. Both algorithms were written by hand
without standard-library sorting functions and tested on five input patterns
(random, ascending, descending, nearly sorted, and with duplicates) with
$n = 1{,}000$, $10{,}000$, and $100{,}000$, using 30 repetitions per condition
and fixed seeds; Bubble Sort was not measured at $n = 100{,}000$ because of its
running cost. Execution time, comparisons, swaps or element moves, and
recursion depth were recorded. For $n = 10{,}000$ and random
data, the median time was 285.4 ms for Bubble Sort and 1.75 ms for Merge Sort,
about 163 times lower. Bubble Sort always performed $n(n-1)/2$ comparisons and
its time varied with the number of swaps; Merge Sort took between 10 and 20 ms
at $n = 100{,}000$, depending on the pattern. The measured growth was consistent with $O(n^2)$ and $O(n\log n)$, and the input pattern changed the times by a factor close to 2 in both algorithms. Merge Sort's advantage came at the cost of $O(n)$ auxiliary memory.}
 
\vspace{0.15cm}
 
\noindent\textbf{Keywords:}
sorting algorithms, Bubble Sort, Merge Sort,
performance, experimental analysis.
 
}
\end{otherlanguage}
 
\vspace{0.5cm}
}]
 
% -------------------- CUERPO DEL ARTICULO --------------------
 
\section{Introducción}
 
Los algoritmos de ordenamiento son procedimientos fundamentales
en el procesamiento de datos, debido a que permiten organizar
elementos de acuerdo con un criterio determinado. Su utilización
facilita la búsqueda, el análisis y la presentación de información
en diferentes aplicaciones informáticas. Sin embargo, los
algoritmos que resuelven el mismo problema pueden presentar
diferencias importantes en su tiempo de ejecución, uso de memoria
y comportamiento frente a distintos tipos de entrada.
 
Entre los algoritmos de ordenamiento se encuentran Bubble Sort
y Merge Sort. Bubble Sort compara elementos adyacentes y los
intercambia cuando se encuentran en un orden incorrecto. Su
implementación es sencilla, aunque su complejidad temporal
promedio y en el peor caso es cuadrática. Por otro lado,
Merge Sort divide los datos en partes más pequeñas, ordena
cada una y combina los resultados. Este procedimiento presenta
una complejidad temporal de $O(n\log n)$ en sus casos mejor,
promedio y peor, aunque necesita memoria auxiliar adicional
para la combinación.
 
El rendimiento de estos algoritmos puede variar según la cantidad
de elementos y el patrón de entrada utilizado. Las versiones de
Bubble Sort con parada anticipada se benefician de las listas que
ya están ordenadas, mientras que la versión canónica realiza el
mismo número de comparaciones sin importar el patrón. Una lista
desordenada de gran tamaño, por su parte, evidencia las diferencias
entre las estrategias de ambos algoritmos. Por esta razón, resulta
conveniente estudiar su comportamiento bajo diferentes condiciones
experimentales.
 
La comparación empírica de algoritmos de ordenamiento ya ha sido
abordada mediante análisis estadístico (Salgado et al., 2022) y
mediante la medición de la eficiencia de implementaciones recursivas
e iterativas (Murillo y Caamaño, 2013), mientras que la complejidad
de los algoritmos clásicos de ordenamiento se encuentra descrita en
textos de referencia (Cormen et al., 2022; Knuth, 1998).
 
El objetivo de este estudio es comparar experimentalmente
Bubble Sort y Merge Sort, implementados en C++, considerando
distintos tamaños y patrones de entrada. Para ello, se registraron
el tiempo de ejecución, el número de comparaciones, los intercambios
o movimientos de elementos y la profundidad máxima de recursión
cuando corresponde.
 
Como expectativa descriptiva, se planteó que Bubble Sort
presentaría un crecimiento del tiempo de ejecución más pronunciado
que Merge Sort al aumentar el tamaño de la entrada, y que Merge Sort
mostraría un crecimiento compatible con su complejidad
$O(n\log n)$. Estas expectativas se contrastan con los datos
obtenidos en la Discusión.
 
\section{Metodología experimental}
 
\subsection{Selección de los algoritmos}
 
Para realizar la comparación experimental se seleccionaron
Bubble Sort y Merge Sort, debido a que utilizan estrategias
diferentes para resolver el problema del ordenamiento. Bubble
Sort se basa en comparaciones e intercambios entre elementos
adyacentes, mientras que Merge Sort emplea la estrategia de
divide y vencerás. Esta selección permite estudiar las diferencias
entre un algoritmo de implementación sencilla y otro que divide
recursivamente el problema en partes más pequeñas. Ambos se
programaron a mano en C++, sin emplear funciones de la biblioteca
estándar para ordenar, copiar o intercambiar elementos.
 
\subsubsection{Bubble Sort}
 
Se utilizó la versión canónica de Bubble Sort, sin condición de
terminación anticipada. En la pasada $i$ se recorre el arreglo de
izquierda a derecha, se compara cada elemento $A[j]$ con el
siguiente y se intercambian cuando $A[j] > A[j+1]$, de modo que el
mayor elemento de la parte no ordenada queda al final.
 
Como no existe parada anticipada, el algoritmo realiza siempre
$n(n-1)/2$ comparaciones, por lo que su complejidad temporal es
$O(n^2)$ en los casos mejor, promedio y peor. Una variante con
parada anticipada tendría $O(n)$ en el mejor caso, pero no se
empleó en este estudio. Su memoria auxiliar es $O(1)$, debido a
que ordena los elementos en el arreglo original.
 
La facilidad de implementación de Bubble Sort permite utilizarlo
como referencia para estudiar el efecto del patrón de entrada
sobre el rendimiento. No obstante, su número de comparaciones
e intercambios aumenta considerablemente con el tamaño de la
entrada.
 
\subsubsection{Merge Sort}
 
Merge Sort se implementó mediante una variante recursiva
descendente basada en la estrategia de divide y vencerás. El
procedimiento divide el arreglo en dos mitades, ordena
recursivamente cada una y combina ambas en una secuencia
ordenada, usando un arreglo auxiliar de $n$ elementos.
 
Su complejidad temporal es $O(n\log n)$ en los casos mejor,
promedio y peor, y su memoria auxiliar es $O(n)$. La profundidad
máxima de recursión se registró como el número de niveles de
llamadas, contando el nivel inicial.
 
Aunque Merge Sort requiere una implementación más elaborada
que Bubble Sort, su estructura permite abordar entradas grandes
con un crecimiento temporal teórico menor. La utilización de
memoria adicional constituye una diferencia importante que
debe tenerse en cuenta al comparar ambos algoritmos.
 
\subsubsection{Complejidad y memoria}
 
La Tabla~\ref{tab:complejidad} presenta las características
teóricas de las variantes utilizadas. Estas propiedades se
emplearon como referencia para interpretar los resultados
experimentales, sin sustituir las mediciones reales. La memoria
corresponde al espacio auxiliar adicional requerido por cada
implementación.
 
\begin{table}[t]
\centering
\caption{Características teóricas de los algoritmos.}\label{tab:complejidad}
\small
\setlength{\tabcolsep}{3pt}
\begin{tabularx}{\columnwidth}{@{}>{\raggedright\arraybackslash}p{2.6cm}>{\raggedright\arraybackslash}X>{\raggedright\arraybackslash}X@{}}
\toprule
 & \textbf{Bubble Sort} & \textbf{Merge Sort} \\
\midrule
Estrategia & Intercambios adyacentes & Divide y vencerás \\
Mejor caso & $O(n^2)$ & $O(n\log n)$ \\
Caso promedio & $O(n^2)$ & $O(n\log n)$ \\
Peor caso & $O(n^2)$ & $O(n\log n)$ \\
Memoria auxiliar & $O(1)$ & $O(n)$ \\
Estabilidad & Estable & Estable \\
Recursión & No & Sí \\
\bottomrule
\end{tabularx}
\end{table}
 
\subsubsection{Estabilidad}
 
La estabilidad es la propiedad que permite conservar el orden
relativo original de los elementos que poseen la misma clave.
Esta característica puede ser importante cuando se ordenan
registros que contienen varios campos.
 
En Bubble Sort solo se intercambian elementos cuando
$A[j] > A[j+1]$, es decir, nunca se intercambian claves iguales.
En Merge Sort, durante la combinación, se toma primero el elemento
de la mitad izquierda cuando ambas claves son iguales
($A[i] \leq A[j]$). Por ello, ambas implementaciones son estables.
Esta propiedad se dedujo del código y no se verificó
experimentalmente, porque los elementos ordenados fueron enteros
simples, que no permiten distinguir claves iguales.
 
\subsection{Patrones de entrada}
 
Se utilizaron cinco patrones de entrada, descritos en la
Tabla~\ref{tab:patrones}, para estudiar no solo el caso de una lista
desordenada al azar, sino también el comportamiento de los
algoritmos cuando los datos ya presentan cierto orden o valores
repetidos.
 
\begin{table}[t]
\centering
\caption{Patrones de entrada utilizados.}\label{tab:patrones}
\small
\setlength{\tabcolsep}{3pt}
\begin{tabularx}{\columnwidth}{@{}>{\raggedright\arraybackslash}p{2.2cm}>{\raggedright\arraybackslash}X@{}}
\toprule
\textbf{Patrón} & \textbf{Descripción} \\
\midrule
Aleatorio & Enteros uniformes entre 1 y 1\,000\,000. \\
Ascendente & $1, 2, \dots, n$. \\
Descendente & $n, n-1, \dots, 1$. \\
Casi ordenado & Arreglo ascendente con $n/40$ intercambios aleatorios (cerca de 5\,\% de los elementos movidos). \\
Con duplicados & Quince valores únicos, entre 1 y 15, repetidos. \\
\bottomrule
\end{tabularx}

\end{table}
 
\subsection{Tamaños, repeticiones y semillas}
 
Se consideraron tres tamaños de entrada,
$n \in \{1\,000;\ 10\,000;\ 100\,000\}$. Cada combinación de
algoritmo, patrón y tamaño se repitió 30 veces, excepto Bubble Sort
con $n = 100\,000$, que se omitió. En ese tamaño, Bubble Sort realiza
cerca de $5\times10^{9}$ comparaciones por ejecución, lo que vuelve
poco práctico repetirlo 30 veces por patrón. Por ello su evaluación
se limitó a $n = 1\,000$ y $n = 10\,000$, mientras que Merge Sort se
evaluó con los tres tamaños. En total se obtuvieron 750 ejecuciones:
300 de Bubble Sort y 450 de Merge Sort.
 
Para facilitar la repetición del experimento se utilizó un generador
congruencial lineal propio, en lugar de \texttt{rand()}, de modo que
los datos son idénticos en cualquier equipo y compilador. La semilla
base fue 12345 y la repetición $r$ usó la semilla $12345 + r$, es
decir, de 12346 a 12375. En cada repetición, ambos algoritmos
recibieron exactamente los mismos datos de entrada.
 
\subsection{Métricas y registro de datos}
 
Se registraron cinco medidas por ejecución: el tiempo de ordenamiento
en milisegundos, medido con \texttt{std::chrono::steady\_clock}; el
número de comparaciones entre elementos; los intercambios, que en
Bubble Sort son intercambios de pares y en Merge Sort son los movimientos
de elementos al arreglo auxiliar, porque este algoritmo no intercambia
como tal; y la profundidad máxima de recursión, que solo se aplica a
Merge Sort y vale 0 en Bubble Sort.
 
El tiempo se midió únicamente durante el ordenamiento: la generación
de los datos, la verificación del resultado y la escritura del archivo
quedaron fuera del intervalo cronometrado. Cada fila del archivo CSV
corresponde a una ejecución e incluye, además, el algoritmo, el
tamaño, el patrón, el número de repetición, la semilla y el
identificador de la versión del código. Después de cada ejecución se
comprobó que el arreglo quedara ordenado y no se detectó ningún error en
las 750 ejecuciones. Los contadores forman parte del trabajo realizado
dentro de los algoritmos medidos, por lo que añaden un pequeño costo al
tiempo, igual para ambos.
 
\subsection{Entorno de ejecución}
 
Las pruebas se realizaron en un equipo con procesador Intel Core
i5-7400 a 3,00 GHz (4 núcleos y 4 hilos), 8 GB de memoria RAM DDR4 a
2133 MHz y sistema operativo Windows 11. El código se compiló con
TDM-GCC 4.9.2 de 64 bits desde Dev-C++, con las opcion
\texttt{-static-libgcc} y sin indicadores de
optimización (\texttt{-O}). Durante la corrida final se cerraron los demás programas. Los
tiempos obtenidos corresponden a este equipo y pueden variar en otro.
 
\subsection{Análisis de los datos}
 
Para cada algoritmo, patrón y tamaño se calcularon la media, la
mediana, la desviación estándar, el mínimo y el máximo del tiempo, y
el promedio de los contadores, con un script de Python (\texttt{pandas}
y \texttt{matplotlib}). El código, los datos y el script se
encuentran en el repositorio del proyecto, y los
resultados corresponden al código del commit \texttt{2c4e2dc}.
 
\section{Resultados}
 
\subsection{Tiempos de ejecución}
 
La Tabla~\ref{tab:tbubble} y la Tabla~\ref{tab:tmerge} resumen los
tiempos de Bubble Sort y de Merge Sort. Con $n = 10\,000$, la mediana de
Bubble Sort estuvo entre 125,5 ms (ascendente) y 286,9 ms (con
duplicados), mientras que la de Merge Sort estuvo entre 0,85 ms
(ascendente) y 1,75 ms (aleatorio). Con datos aleatorios, la mediana de
Bubble Sort fue unas 163 veces mayor que la de Merge Sort. Con
$n = 1\,000$ esa razón fue de unas 17 veces.
 
Con $n = 100\,000$ solo se midió Merge Sort, cuyas medianas
estuvieron entre 10,4 ms (ascendente) y 20,0 ms (aleatorio).
 
 
\begin{table}[t]
\centering
\caption{Estadística descriptiva del tiempo de ejecución (ms) de Bubble Sort, con 30 repeticiones por condición.}\label{tab:tbubble}
\footnotesize
\setlength{\tabcolsep}{3pt}
\fitcol{\begin{tabular}{llrrrrr}
\toprule
Patrón & $n$ & Media & Mediana & Desv.\ est. & Mín. & Máx. \\
\midrule
Aleatorio & 1\,000 & 2,21 & 2,13 & 0,17 & 2,06 & 2,83 \\
Aleatorio & 10\,000 & 295,97 & 285,38 & 24,92 & 278,61 & 404,05 \\
\addlinespace[2pt]
Ascendente & 1\,000 & 1,26 & 1,23 & 0,09 & 1,18 & 1,62 \\
Ascendente & 10\,000 & 126,56 & 125,46 & 2,87 & 123,92 & 133,92 \\
\addlinespace[2pt]
Descendente & 1\,000 & 2,17 & 2,13 & 0,09 & 2,07 & 2,42 \\
Descendente & 10\,000 & 235,82 & 222,02 & 55,57 & 217,95 & 524,37 \\
\addlinespace[2pt]
Casi ordenado & 1\,000 & 1,51 & 1,38 & 0,42 & 1,32 & 3,54 \\
Casi ordenado & 10\,000 & 156,48 & 147,69 & 19,35 & 140,25 & 207,74 \\
\addlinespace[2pt]
Con duplicados & 1\,000 & 2,13 & 2,05 & 0,19 & 2,01 & 2,79 \\
Con duplicados & 10\,000 & 300,35 & 286,85 & 28,04 & 276,77 & 393,21 \\
\bottomrule
\end{tabular}}
 
\end{table}
 
 
\begin{table}[t]
\centering
\caption{Estadística descriptiva del tiempo de ejecución (ms) de Merge Sort, con 30 repeticiones por condición.}\label{tab:tmerge}
\footnotesize
\setlength{\tabcolsep}{3pt}
\fitcol{\begin{tabular}{llrrrrr}
\toprule
Patrón & $n$ & Media & Mediana & Desv.\ est. & Mín. & Máx. \\
\midrule
Aleatorio & 1\,000 & 0,166 & 0,128 & 0,139 & 0,116 & 0,868 \\
Aleatorio & 10\,000 & 1,869 & 1,754 & 0,322 & 1,628 & 3,195 \\
Aleatorio & 100\,000 & 20,211 & 19,975 & 0,603 & 19,708 & 22,615 \\
\addlinespace[2pt]
Ascendente & 1\,000 & 0,069 & 0,061 & 0,017 & 0,060 & 0,119 \\
Ascendente & 10\,000 & 0,880 & 0,848 & 0,094 & 0,792 & 1,223 \\
Ascendente & 100\,000 & 10,642 & 10,400 & 0,636 & 10,167 & 13,072 \\
\addlinespace[2pt]
Descendente & 1\,000 & 0,106 & 0,074 & 0,105 & 0,061 & 0,589 \\
Descendente & 10\,000 & 0,939 & 0,880 & 0,144 & 0,805 & 1,365 \\
Descendente & 100\,000 & 10,935 & 10,479 & 1,544 & 10,132 & 18,408 \\
\addlinespace[2pt]
Casi ordenado & 1\,000 & 0,069 & 0,066 & 0,009 & 0,065 & 0,112 \\
Casi ordenado & 10\,000 & 0,968 & 0,930 & 0,103 & 0,880 & 1,335 \\
Casi ordenado & 100\,000 & 11,363 & 11,306 & 0,264 & 10,996 & 12,031 \\
\addlinespace[2pt]
Con duplicados & 1\,000 & 0,110 & 0,101 & 0,029 & 0,099 & 0,253 \\
Con duplicados & 10\,000 & 1,275 & 1,240 & 0,123 & 1,189 & 1,811 \\
Con duplicados & 100\,000 & 15,251 & 14,743 & 1,520 & 14,485 & 22,694 \\
\bottomrule
\end{tabular}}
 
\end{table}
 
 
\subsection{Crecimiento con el tamaño y variabilidad}
 
La Figura~\ref{fig:n-tiempo} muestra el tiempo mediano en función de
$n$. Al pasar de $n = 1\,000$ a $n = 10\,000$, la mediana de Bubble Sort
se multiplicó entre 102 y 140 veces, según el patrón. En Merge Sort,
el mismo cambio multiplicó la mediana entre 11,8 y 14,2 veces, y el paso
de $10\,000$ a $100\,000$ la multiplicó entre 11,4 y 12,3 veces.
 
La Figura~\ref{fig:boxplot} presenta la distribución de los 30 tiempos
con $n = 10\,000$. En Bubble Sort, el coeficiente de variación estuvo
entre 2,3\,\% y 23,6\,\% con $n = 10\,000$. En Merge Sort fue mayor con
$n = 1\,000$ (de 13,2\,\% a 98,6\,\%) y menor con $n = 100\,000$ (de
2,3\,\% a 14,1\,\%). También se observaron valores atípicos aislados,
como un máximo de 524 ms en Bubble Sort descendente con $n = 10\,000$,
frente a una mediana de 222 ms.
 
 
\begin{figure}[t]
\centering
\includegraphics[width=\columnwidth]{fig1_n_vs_tiempo.png}
\caption{Tiempo mediano de ejecución en función de $n$, en escala logarítmica en ambos ejes.}\label{fig:n-tiempo}
\end{figure}
 
\begin{figure}[t]
\centering
\includegraphics[width=\columnwidth]{fig2_boxplot.png}
\caption{Distribución de los tiempos de ejecución con $n = 10\,000$ (escala logarítmica).}\label{fig:boxplot}
\end{figure}
 
 
\subsection{Comparaciones, intercambios y profundidad}
 
La Tabla~\ref{tab:cbubble} muestra que Bubble Sort realizó siempre
$n(n-1)/2$ comparaciones: 499\,500 con $n = 1\,000$ y 49\,995\,000 con
$n = 10\,000$, sin importar el patrón. Lo que cambió con el patrón fue el
número de intercambios: 0 en el patrón ascendente, $n(n-1)/2$ en el
descendente y cerca de la mitad de esa cantidad en los patrones
aleatorio y con duplicados.
 
En Merge Sort (Tabla~\ref{tab:cmerge}) los movimientos al arreglo
auxiliar fueron iguales en todos los patrones para cada tamaño (9\,976,
133\,616 y 1\,668\,928), mientras que las comparaciones dependieron del
patrón: con $n = 10\,000$ fueron 120\,457 en el aleatorio y 64\,608 en el
descendente. La profundidad de recursión fue 11, 15 y 18 para
$n = 1\,000$, $10\,000$ y $100\,000$. La Figura~\ref{fig:contadores}
resume estos contadores con $n = 10\,000$.
 
 
\begin{table}[t]
\centering
\caption{Contadores de Bubble Sort (promedio de 30 repeticiones). La profundidad de recursión es 0.}\label{tab:cbubble}
\footnotesize
\setlength{\tabcolsep}{3pt}
\fitcol{\begin{tabular}{llrr}
\toprule
Patrón & $n$ & Comparaciones & Intercambios \\
\midrule
Aleatorio & 1\,000 & 499\,500 & 249\,025 \\
Aleatorio & 10\,000 & 49\,995\,000 & 25\,015\,115 \\
\addlinespace[2pt]
Ascendente & 1\,000 & 499\,500 & 0 \\
Ascendente & 10\,000 & 49\,995\,000 & 0 \\
\addlinespace[2pt]
Descendente & 1\,000 & 499\,500 & 499\,500 \\
Descendente & 10\,000 & 49\,995\,000 & 49\,995\,000 \\
\addlinespace[2pt]
Casi ordenado & 1\,000 & 499\,500 & 16\,656 \\
Casi ordenado & 10\,000 & 49\,995\,000 & 1\,623\,316 \\
\addlinespace[2pt]
Con duplicados & 1\,000 & 499\,500 & 233\,856 \\
Con duplicados & 10\,000 & 49\,995\,000 & 23\,363\,148 \\
\bottomrule
\end{tabular}}
 
\end{table}
 
 
\begin{table}[t]
\centering
\caption{Contadores de Merge Sort (promedio de 30 repeticiones). Los movimientos son copias al arreglo auxiliar.}\label{tab:cmerge}
\footnotesize
\setlength{\tabcolsep}{3pt}
\fitcol{\begin{tabular}{llrrr}
\toprule
Patrón & $n$ & Comparaciones & Movimientos & Profundidad \\
\midrule
Aleatorio & 1\,000 & 8\,712 & 9\,976 & 11 \\
Aleatorio & 10\,000 & 120\,457 & 133\,616 & 15 \\
Aleatorio & 100\,000 & 1\,536\,397 & 1\,668\,928 & 18 \\
\addlinespace[2pt]
Ascendente & 1\,000 & 5\,044 & 9\,976 & 11 \\
Ascendente & 10\,000 & 69\,008 & 133\,616 & 15 \\
Ascendente & 100\,000 & 853\,904 & 1\,668\,928 & 18 \\
\addlinespace[2pt]
Descendente & 1\,000 & 4\,932 & 9\,976 & 11 \\
Descendente & 10\,000 & 64\,608 & 133\,616 & 15 \\
Descendente & 100\,000 & 815\,024 & 1\,668\,928 & 18 \\
\addlinespace[2pt]
Casi ordenado & 1\,000 & 7\,163 & 9\,976 & 11 \\
Casi ordenado & 10\,000 & 107\,131 & 133\,616 & 15 \\
Casi ordenado & 100\,000 & 1\,403\,802 & 1\,668\,928 & 18 \\
\addlinespace[2pt]
Con duplicados & 1\,000 & 8\,588 & 9\,976 & 11 \\
Con duplicados & 10\,000 & 118\,268 & 133\,616 & 15 \\
Con duplicados & 100\,000 & 1\,502\,989 & 1\,668\,928 & 18 \\
\bottomrule
\end{tabular}}
 
\end{table}
 
 
\begin{figure}[t]
\centering
\includegraphics[width=\columnwidth]{fig3_contadores.png}
\caption{Intercambios de Bubble Sort y comparaciones de Merge Sort con $n = 10\,000$ (promedio de 30 repeticiones). La línea discontinua marca las comparaciones de Bubble Sort.}\label{fig:contadores}
\end{figure}
 
 
\section{Discusión}
 
\subsection{Relación con la complejidad teórica}
 
Los crecimientos medidos son compatibles con las complejidades
teóricas. Para Bubble Sort, el número de comparaciones coincide
exactamente con $n(n-1)/2$ y, al multiplicar $n$ por 10, la teoría
cuadrática predice un tiempo unas 100 veces mayor; las medianas
crecieron entre 102 y 140 veces, un poco más en los patrones aleatorio
y con duplicados. Para Merge Sort, la teoría $O(n\log n)$ predice un
factor de 13,3 al pasar de $1\,000$ a $10\,000$ y de 12,5 al pasar de
$10\,000$ a $100\,000$, y los valores medidos (11,8 a 14,2 y 11,4 a
12,3) estuvieron cerca. La expectativa descriptiva planteada en la
introducción se cumplió: el crecimiento de Bubble Sort fue más
pronunciado y el de Merge Sort fue compatible con $O(n\log n)$.
 
\subsection{Efecto del patrón de entrada}
 
En Bubble Sort, las comparaciones fueron fijas, por lo que las diferencias
entre patrones provienen de los intercambios y del comportamiento del
procesador. El patrón ascendente no realiza intercambios y fue el más
rápido (125,5 ms con $n = 10\,000$). Sin embargo, el patrón descendente,
con el doble de intercambios que el aleatorio, fue más rápido que este
(222,0 ms frente a 285,4 ms). Esto sugiere que, además del número de
intercambios, influye la regularidad de las comparaciones, por ejemplo
la predicción de saltos del procesador; esta explicación no se verificó
con mediciones adicionales.
 
En Merge Sort, los patrones ascendente, descendente y casi ordenado
requirieron menos comparaciones (entre 64\,608 y 107\,131 con
$n = 10\,000$) que el aleatorio (120\,457) y tardaron menos (de 0,85 a
0,93 ms frente a 1,75 ms). El patrón con duplicados tuvo un número
de comparaciones parecido al aleatorio (118\,268) pero un tiempo menor
(1,24 ms), lo que de nuevo indica que el tiempo no depende solo del
número de comparaciones.
 
\subsection{Duplicados y estabilidad}
 
Como ambas implementaciones son estables, los valores repetidos no
alteran su orden relativo. En términos de rendimiento, con duplicados
Bubble Sort hizo algo menos de intercambios que con datos aleatorios
(23,4 millones frente a 25,0 millones con $n = 10\,000$), porque no
intercambia claves iguales, y Merge Sort hizo menos comparaciones
(118\,268 frente a 120\,457).
 
\subsection{Memoria}
 
Merge Sort necesitó un arreglo auxiliar del mismo tamaño que la entrada,
que con $n = 100\,000$ enteros de 4 bytes equivale a unos 0,4 MB, mientras
que Bubble Sort ordenó en el sitio. La ventaja de tiempo de Merge Sort
se obtuvo, por tanto, a cambio de memoria adicional. La memoria no se
midió experimentalmente; se consideró el valor teórico.
 
\subsection{Limitaciones y amenazas a la validez}
 
Los resultados provienen de un solo equipo, un solo compilador con su
configuración por defecto y datos sintéticos de enteros, por lo que no
deben interpretarse como tiempos universales. Bubble Sort no se midió con
$n = 100\,000$, de modo que la comparación directa entre los dos
algoritmos se limita a $n = 1\,000$ y $n = 10\,000$. Con $n = 1\,000$,
Merge Sort tarda unas décimas de milisegundo y su variabilidad relativa
fue alta (hasta 98,6\,\%), por lo que se emplearon las medianas. Los
contadores añaden un costo a los tiempos y la estabilidad se dedujo del
código sin verificarse experimentalmente. Por último, los valores
atípicos pueden deberse a la actividad del sistema operativo.
 
\section{Conclusiones}
 
\begin{enumerate}
\setlength{\itemsep}{2pt}
\item Con $n = 10\,000$ y datos aleatorios, la mediana fue de 285,4 ms en
Bubble Sort y de 1,75 ms en Merge Sort, unas 163 veces menor, y con
$n = 100\,000$ Merge Sort tardó entre 10 y 20 ms.
\item Bubble Sort realizó siempre $n(n-1)/2$ comparaciones, mientras que
Merge Sort realizó entre 64\,608 y 120\,457 con $n = 10\,000$, según el
patrón.
\item El crecimiento medido fue compatible con $O(n^2)$ en Bubble Sort y
con $O(n\log n)$ en Merge Sort.
\item El patrón de entrada modificó los tiempos en un factor cercano a 2
en ambos algoritmos, y esa variación no se explica solo por el número de
comparaciones o intercambios.
\item La ventaja de tiempo de Merge Sort requiere un arreglo auxiliar de
orden $O(n)$, mientras que Bubble Sort usa memoria auxiliar constante.
\end{enumerate}
 
\section*{Referencias}
 
{\small\sloppy
\setlength{\parindent}{0pt}
\setlength{\parskip}{3pt}
 
\hangindent=0.5cm\hangafter=1
Cormen, T. H., Leiserson, C. E., Rivest, R. L., \& Stein, C. (2022).
\textit{Introduction to algorithms} (4.ª ed.). MIT Press.\par
 
\hangindent=0.5cm\hangafter=1
Knuth, D. E. (1998). \textit{The art of computer programming: Vol. 3.
Sorting and searching} (2.ª ed.). Addison-Wesley.\par
 
\hangindent=0.5cm\hangafter=1
Murillo, J., \& Caamaño, S. (2013). Comparación entre algoritmos
recursivos e iterativos y su medición en términos de eficiencia.
\textit{Uniciencia, 27}(1), 341--350.
\url{https://www.redalyc.org/pdf/4759/475947762020.pdf}\par
 
\hangindent=0.5cm\hangafter=1
Salgado, M., Pérez, M., \& Albarrán, S. (2022). Algoritmos de ordenamiento,
un análisis estadístico comparativo. \textit{Ciencia e Ingeniería,
43}(3), 273--280. \url{https://doi.org/10.53766/cei/2022.43.03.05}\par
}
 
 
\end{document}
 